\chapter*{Lista de Símbolos}
\addcontentsline{toc}{chapter}{Lista de Símbolos}

\begin{tabularx}{\textwidth}{>{\raggedright\arraybackslash}p{3.0cm} >{\raggedright\arraybackslash}X}
$t$ & Tempo da amostra, em segundos\\
$t_{us}$ & Timestamp registrado pelo microcontrolador, em microssegundos\\
$\Delta t_i$ & Intervalo entre as amostras $i-1$ e $i$\\
$f_s$ & Frequência de amostragem estimada a partir de $\Delta t$\\

$\mathbf{p}$ & Posição ou coordenada de um ponto\\
$\mathbf{p}_{drone,G}$ & Posição do drone no referencial global\\
$\mathbf{v}$ & Velocidade do drone no referencial global\\
$\mathbf{a}$ & Aceleração linear\\
$\tilde{\mathbf{a}}_B$ & Aceleração específica medida pela IMU no referencial do corpo\\
$\mathbf{a}_G$ & Aceleração linear no referencial global após compensação da gravidade\\
$\mathbf{g}_G$ & Vetor gravidade no referencial global\\
$g$ & Módulo da aceleração da gravidade, adotado como $9{,}81~\text{m/s}^2$\\

$\boldsymbol{\omega}$ & Velocidade angular medida pelo giroscópio\\
$\mathbf{b}_a,\mathbf{b}_\omega$ & Viés do acelerômetro e do giroscópio\\
$\boldsymbol{\eta}_a,\boldsymbol{\eta}_\omega$ & Ruído de medição do acelerômetro e do giroscópio\\
$d$ & Deriva acumulada pela integração temporal dos erros dos sensores\\

$\phi$ & Ângulo de \textit{roll}, rotação em torno do eixo $X$\\
$\theta$ & Ângulo de \textit{pitch}, rotação em torno do eixo $Y$\\
$\psi$ & Ângulo de \textit{yaw}, rotação em torno do eixo $Z$\\
$\mathbf{q}$ & Quaternion de atitude no formato $[q_w,q_x,q_y,q_z]$\\
$\beta$ & Ganho do filtro de Madgwick\\
$\mathbf{R}_{B}^{G}$ & Matriz de rotação que transforma vetores do corpo do drone para o referencial global\\
$R_x,R_y,R_z$ & Matrizes elementares de rotação em torno dos eixos $X$, $Y$ e $Z$\\

$\Delta x_f,\Delta y_f$ & Deslocamentos incrementais do sensor óptico de fluxo, em centímetros no CSV\\
$h$ & Altura medida pelo ToF\\
$h_0$ & Altura nominal usada na escala do fluxo óptico\\
$k_h$ & Fator de escala do ToF, dado por $h/h_0$\\
$\mathbf{p}_{flow}$ & Posição acumulada a partir do fluxo óptico no plano $XY$\\
$\mathbf{p}_{fused}$ & Posição final após combinação de fluxo óptico com propagação inercial nos intervalos sem eventos\\

$\alpha_L$ & Ângulo de varredura do LiDAR, em centésimos de grau no CSV bruto\\
$\theta_L$ & Ângulo de varredura do LiDAR convertido para radianos\\
$r_L$ & Distância medida pelo LiDAR\\
$I_L$ & Intensidade do retorno LiDAR\\
$\gamma_Y$ & Ângulo físico de varredura/rotação do LiDAR principal em torno do eixo $Y$ do drone\\
$\mathbf{p}_L$ & Coordenada do ponto no referencial local do LiDAR\\
$\mathbf{p}_B$ & Coordenada do ponto no referencial do corpo do drone\\
$\mathbf{p}_G$ & Coordenada global do ponto da nuvem\\
$\mathbf{R}_{B}^{L}$ & Rotação extrínseca entre LiDAR e corpo do drone\\
$\mathbf{t}_{B}^{L}$ & Translação extrínseca entre LiDAR e corpo do drone\\

$Q_1,Q_3$ & Primeiro e terceiro quartis usados na cerca de Tukey\\
$IQR$ & Intervalo interquartil, $IQR=Q_3-Q_1$\\
$\tau_I,\tau_r$ & Limiares automáticos de intensidade e distância calculados por Otsu\\
\end{tabularx}
